\section{Merit order and price spikes --- why price is hard}

\subsection{The merit order}

Electricity is dispatched in cost order. The system operator calls the
cheapest plant first, then the next, until supply meets demand. The
\emph{marginal} plant --- the last one needed --- sets the price for
everyone. This is the merit order.

In practice (Poland, any hour):

\begin{enumerate}
  \item Nuclear and run-of-river hydro: near-zero marginal cost, always on.
  \item Wind and solar: zero fuel cost. They run whenever they produce.
  \item Lignite (r\k{e}gowe kopalnia): cheapest thermal, base-load.
  \item Hard coal: mid-merit.
  \item Gas (CCGT): most expensive common thermal.
  \item Imports, peak hydro, demand response: sets the price only when
        cheaper options are exhausted.
\end{enumerate}

\textbf{Key consequence.} More wind or solar pushes coal/gas down the stack.
When enough renewables enter, the marginal unit flips from gas to coal,
then to lignite. Price falls. If renewables exceed demand, the price can
go \emph{negative}: inflexible generators pay to stay on.

\textbf{Measured in our data.} Wind $>4\,\mathrm{GW}$ (Poland): average price
$87.8\,\mathrm{EUR/MWh}$ vs $111.0\,\mathrm{EUR/MWh}$ at low wind --- a 21\%
discount, recoverable by a single DuckDB query (notebook cell 4).

\subsection{Price formation through a day}

A typical summer day (Poland):

\begin{itemize}
  \item 00:00--06:00: low demand, base-load only. Price $\sim90$ EUR.
  \item 06:00--08:00: morning ramp, industrial demand. Price spikes to
        $130+$ EUR (gas enters).
  \item 09:00--14:00: solar peak. Price falls, sometimes to negative
        (April 2026 had 87 negative hours). LEAR sees this as the
        ``solar trough''.
  \item 15:00--21:00: solar wanes, demand holds high. Evening ramp.
        Gas is marginal again. Price returns to 130--165 EUR.
  \item 22:00--23:00: demand drops; night settle. Prices fall back.
\end{itemize}

This shape appears in our SQL output (notebook cell 3): weekday midday
48 EUR (solar), evening 165 EUR (gas). Weekend midday 48 EUR too, but
evening ramp is smaller (less industrial demand).

\subsection{Why spikes are hard to forecast}

A price spike ($>250\,\mathrm{EUR/MWh}$) happens when:

\begin{itemize}
  \item Unexpected demand peak (cold snap, hot spell) AND
  \item Unexpected supply loss (plant outage, calm wind) AND
  \item Low import capacity (interconnectors congested).
\end{itemize}

All three must align. Our feature set captures the first two partially
(temperature forecast, wind forecast) but misses:
\begin{itemize}
  \item \textbf{Unit outages}: large generators file Urgent Market Messages
        (UMM) on ENTSO-E. We built the feature; it was FLAT in our test.
        The aggregate signal is too coarse and revises after events.
  \item \textbf{Interconnector capacity}: cross-border transfer limits
        change daily. Not yet in our feature set.
  \item \textbf{Gas-day prices}: TTF intraday moves during the forecast window.
        Our proxy (overnight settlement) is a lag.
\end{itemize}

\textbf{Spike MAE benchmark}: LGBM 60.6 EUR/MWh vs naive 77.6 EUR/MWh.
A 22\% improvement, but still 3$\times$ the pooled-hour MAE. This is honest.
Any price forecaster who claims to ``solve'' spikes is selling something.

\subsection{Gas and CO2 as level drivers}

Gas (TTF) sets the variable cost of the marginal CCGT. When TTF is high,
the merit-order price for every gas-dispatched hour rises.

CO2 (EUA permit price) adds to the marginal cost of fossil plants:
lignite has higher specific emissions than gas (more CO2 per MWh), so
high EUA prices narrow the gas-coal spread and can flip merit-order
ranking.

\textbf{In our fuel test (2026-07-17)}: including a TTF index and an
EUA-proxy ETF reduced LEAR winter-2024/25 bias from $-15.9$ to $-4.6$
EUR/MWh ($+2.5$ MAE). The gain concentrated exactly in high-gas months.
Calm-gas period (2025/26): no effect. Merit-order mechanism, measured.

\subsection{Negative prices}

Negative prices arise when wind+solar generation exceeds demand and
inflexible thermal plants cannot ramp down fast enough. They ``pay'' the
grid to take their power so they can keep running (start-up costs exceed
one hour of negative revenue).

Poland negative hours by month (from notebook):
April 2026: 87 hours (peak solar, mild spring demand).
December 2025: 0 hours (low solar, high heating demand).

Negative price hours cannot be handled by MAPE (division by price $\to$
sign change). We use rMAE (MAE relative to naive) throughout.

\subsection{Interview lines}

``The merit order is why solar is our \#1 feature. A large wind ramp 90
minutes before delivery changes the marginal unit from gas to coal --- 30
EUR/MWh difference. Our SHAP and group ablation both recovered this.''

``Spikes need information we don't have: real-time outages, intraday gas
prices, interconnector limits. We tried the outage feature; it was flat.
Admitting that is more credible than claiming to solve it.''

``We use rMAE, not MAPE, because 798 hours in our 3-year dataset had
negative prices. MAPE would give meaningless values or divide-by-zero.''
